\documentclass[sn-basic,Numbered]{sn-jnl}
\usepackage{graphicx}
\usepackage{multirow}
\usepackage{amsmath,amssymb,amsfonts}
\usepackage{booktabs}
\usepackage[T1]{fontenc}
\usepackage{longtable}
\usepackage{array}
\usepackage{tabularx}
\usepackage{xurl}
\urlstyle{same}
% Float placement. Two competing failures. Left at the defaults, with every float
% emitted after the body, all eight figures were deferred to the last two pages;
% opened up far enough to fix that, three pages ended up carrying three floats
% each while nineteen carried none. These are the middle: at most two floats to a
% page, and at least 18 % of any page with floats given to text.
\setcounter{topnumber}{2}
\setcounter{bottomnumber}{1}
\setcounter{totalnumber}{2}
\setcounter{dbltopnumber}{2}
\renewcommand{\topfraction}{0.75}
\renewcommand{\bottomfraction}{0.35}
\renewcommand{\textfraction}{0.18}
\renewcommand{\dbltopfraction}{0.75}
% A float page is only worth making when it is nearly full. Left at 0.65 LaTeX
% collected tables and figures onto pages of their own with no running text and
% large gaps between them, which is the opposite of spreading them out.
\renewcommand{\floatpagefraction}{0.9}
\renewcommand{\dblfloatpagefraction}{0.9}
% And if one is made anyway, pack it from the top rather than distributing the
% floats down the page with stretch between them.
\makeatletter
\setlength{\@fptop}{0pt}
\setlength{\@fpsep}{10pt plus 1fil}
\setlength{\@fpbot}{0pt plus 1fil}
\setlength{\@dblfptop}{0pt}
\setlength{\@dblfpsep}{10pt plus 1fil}
\setlength{\@dblfpbot}{0pt plus 1fil}
\makeatother
% verbatim in a two-column layout has one column to live in; the default
% \small is too wide for a 58-character line and ran 242 pt past the column.
\makeatletter\def\verbatim@font{\ttfamily\scriptsize}\makeatother
\usepackage{calc}
\usepackage{listings}
% T1/Latin Modern silently drops Greek and several maths symbols under XeTeX,
% which matters here because the reliability signals are named by Greek letters.
\usepackage{newunicodechar}
\newunicodechar{α}{\ensuremath{\alpha}}
\newunicodechar{ε}{\ensuremath{\varepsilon}}
\newunicodechar{θ}{\ensuremath{\theta}}
\newunicodechar{ρ}{\ensuremath{\rho}}
\newunicodechar{φ}{\ensuremath{\varphi}}
\newunicodechar{σ}{\ensuremath{\sigma}}
\newunicodechar{τ}{\ensuremath{\tau}}
\newunicodechar{κ}{\ensuremath{\kappa}}
\newunicodechar{Δ}{\ensuremath{\Delta}}
\newunicodechar{−}{\ensuremath{-}}
\newunicodechar{×}{\ensuremath{\times}}
\newunicodechar{±}{\ensuremath{\pm}}
\newunicodechar{≥}{\ensuremath{\geq}}
\newunicodechar{≤}{\ensuremath{\leq}}
\newunicodechar{≈}{\ensuremath{\approx}}
\newunicodechar{→}{\ensuremath{\rightarrow}}
\newunicodechar{⟺}{\ensuremath{\Longleftrightarrow}}
\newunicodechar{⁻}{\ensuremath{^{-}}}
\newunicodechar{¹}{\ensuremath{^{1}}}
\newunicodechar{·}{\ensuremath{\cdot}}
\newunicodechar{§}{\S}
\newunicodechar{…}{\dots}
\newunicodechar{ć}{\'{c}}
\newunicodechar{š}{\v{s}}
\lstset{basicstyle=\ttfamily\scriptsize,breaklines=true,frame=none,columns=fullflexible}
\providecommand{\tightlist}{\setlength{\itemsep}{0pt}\setlength{\parskip}{0pt}}
% sn-jnl loads breakurl, which emits dvips PostScript specials and breaks under
% XeTeX/tectonic. Neutralise it; URLs still typeset, they just do not line-break.
\makeatletter
\providecommand\headerps@out[1]{}
\let\headerps@out\@gobble
% breakurl also replaces \url with a dvips-only implementation (\pdf@box);
% hyperref's \nolinkurl typesets the same text and works under XeTeX.
\AtBeginDocument{\let\url\nolinkurl}
\makeatother
% sn-jnl's \orcidlogo hard-codes Orcidlogo.eps, which XeTeX cannot read.
% Point it at the PDF shipped alongside; the mark renders identically.
\makeatletter
\def\orcidlogo{\raisebox{-0.5pt}{\includegraphics[height=7.5pt]{Orcidlogo-eps-converted-to.pdf}}}
\makeatother
\theoremstyle{thmstyleone}
\newtheorem{theorem}{Theorem}

\begin{document}
\title{Supplementary Material}
\author[1]{\fnm{Yifu}\sur{Zhao}\orcid{https://orcid.org/0009-0004-2363-9269}}\email{p2523269@mpu.edu.mo}
\author[2]{\fnm{Xiaofan}\sur{Zou}\orcid{https://orcid.org/0009-0005-5995-3150}}\email{xiaofanz@shu.edu.cn}
\author[1]{\fnm{Yanxiao}\sur{Li}\orcid{https://orcid.org/0009-0008-3389-1619}}\email{p2525981@mpu.edu.mo}
\author[1]{\fnm{Junhao}\sur{Wei}\orcid{https://orcid.org/0009-0006-0553-2032}}\email{p2312195@mpu.edu.mo}
\author[3]{\fnm{Sio-Kei}\sur{Im}\orcid{https://orcid.org/0000-0002-5599-4300}}\email{marcusim@mpu.edu.mo}
\author*[1]{\fnm{Yapeng}\sur{Wang}\orcid{https://orcid.org/0000-0002-1085-5091}}\email{yapengwang@mpu.edu.mo}
\author[1]{\fnm{Xu}\sur{Yang}\orcid{https://orcid.org/0000-0002-7037-3609}}\email{xuyang@mpu.edu.mo}
\affil*[1]{\orgdiv{Faculty of Applied Sciences}, \orgname{Macao Polytechnic University}, \city{Macao}, \postcode{999078}, \country{China}}
\affil[2]{\orgdiv{School of Mechanical and Electrical Engineering and Automation}, \orgname{Shanghai University}, \city{Shanghai}, \postcode{200444}, \country{China}}
\affil[3]{\orgname{Macao Polytechnic University}, \city{Macao}, \postcode{999078}, \country{China}}
\maketitle
\raggedbottom
\begin{center}\textit{International Journal of Computer Vision}\end{center}
\renewcommand{\thesection}{S\arabic{section}}
\renewcommand{\thetable}{S\arabic{table}}
\emph{Camera-Motion Compensation in Tracking-by-Detection: Accuracy Headroom and Depth-Aware Shared Warps} The released measurements and evaluation summaries provide the numerical sources for the supplementary tables and main-text results.

\section{Reliability-Signal Subsets on MOT17}\label{reliability-signal-subsets-on-mot17}

Section 4.4 identifies the median transfer residual as a compact reliability signal: its AUC of 0.866 matches the best two-signal subset. Table S1 gives the best subset at each size under leave-one-sequence-out over the seven MOT17 sequences.


{\def\LTcaptype{none} % do not increment counter
\begin{table*}[tp]
\centering
\caption{Reliability-signal subsets and held-out AUC.}\label{tab:1}
\fontsize{8.5}{10.5}\selectfont
\setlength{\tabcolsep}{3pt}
\renewcommand{\arraystretch}{1.12}
\begin{tabular*}{\textwidth}{@{\extracolsep{\fill}}llll@{}}
\toprule
Subset size & Best subset & Held-out AUC & Worst sequence \\
\midrule

\bottomrule

\textbf{1} & \textbf{ε} & \textbf{0.866} & 0.779 \\
2 & n + ε & 0.866 & 0.779 \\
3 & ρ + n + ε & 0.864 & 0.780 \\
4 & ρ + n + ε + κ & 0.860 & 0.766 \\
5 & ρ + n + ε + κ + φ & 0.854 & 0.770 \\
6 & all & 0.774 & 0.614 \\
\end{tabular*}
\end{table*}

}

Figure 1 plots every subset. The release includes measurements for all 63 non-empty subsets.

\section{Measurement Conventions}\label{measurement-conventions}

The target-fitted similarity uses least squares over camera-induced correspondences derived from sensor ego-motion and annotated object depths. Background-fitted models use robust estimation over monocular-depth correspondences. The per-target reference uses annotated association; the global homography uses estimated background depth on moving frames and the annotation-anchored near-static fallback defined in the main text. All reported MOT17 oracle comparisons use strict reference substitution.


{\def\LTcaptype{none} % do not increment counter
\begin{table}[tp]
\centering
\caption{Within-frame residual spread after target-fitted similarity over 4,318 moving frames.}\label{tab:2}
\fontsize{8.5}{10.5}\selectfont
\setlength{\tabcolsep}{3pt}
\renewcommand{\arraystretch}{1.12}
\begin{tabularx}{\columnwidth}{@{}>{\raggedright\arraybackslash}Xllll@{}}
\toprule
Estimator & Median & p90 & Max & \textgreater{} 5 px \\
\midrule

\bottomrule

Least squares & \textbf{2.431 px} & \textbf{10.114} & \textbf{67.7} & \textbf{27.12 \%} \\
\end{tabularx}
\end{table}

}

\section{Controls for Class-Specific Outcomes}\label{controls-for-class-specific-outcomes}

Both are measured over the 31,470 object-frames that lie in moving frames with at least two annotated objects --- 23,348 car and 8,122 pedestrian or cyclist. \textbf{Depth-Offset Control.} The median of \textbar log(z\_object / z\_frame-median)\textbar{} is \textbf{0.325 for cars} and \textbf{0.246 for pedestrians and cyclists}; a one-sided Mann--Whitney test for larger pedestrian/cyclist depth offsets returns p = 1.000. These measurements separate depth offset from the class-level tracking outcome. \textbf{Width-Normalised Exposure.} The disagreement between an object's own warp and the shared warp exceeds one third of the object's box width --- roughly where a pure translation drops IoU below the gate --- in \textbf{3.28 \%} of car object-frames and \textbf{3.23 \%} of pedestrian ones. The two classes have similar width-normalised exposure. These controls distinguish class-level tracking outcomes from depth offset and width-normalised exposure alone.

\end{document}